\documentclass[11pt,a4paper]{article}
\usepackage[margin=1in]{geometry}
\usepackage{amsmath,amssymb,amsthm}
\usepackage{algorithm}
\usepackage{algpseudocode}
\usepackage{booktabs}
\usepackage{hyperref}

\theoremstyle{definition}
\newtheorem{definition}{\textbf{Definition}}[section]
\newtheorem{theorem}{\textbf{Theorem}}[section]
\newtheorem{lemma}[theorem]{\textbf{Lemma}}
\newtheorem{remark}{\textbf{Remark}}[section]

\title{\textbf{Deep Interest Network and Target-Dependent Attention (DIN)}}
\author{\textbf{Team Scaibu} \\ \small Email: \texttt{jaydeep.9664920749@gmail.com} \\ \small \textbf{Architecture and Core Engine Team}}
\date{\textbf{October 2026}}

\begin{document}
\maketitle

\section{Definitions and Cognitive Mental Models}

\subsection{Formal Mathematical Definition}
\textbf{Definition (Target-Dependent Attention, Multi-Faceted User Interest, and Activation Units):}
Let $\mathbf{v}_A \in \mathbb{R}^D$ denote the embedding vector of a candidate ad or item to be scored, and let $\mathcal{H} = \{\mathbf{e}_1, \mathbf{e}_2, \dots, \mathbf{e}_H\}$ denote a user's historical behavior sequence of $H$ previously interacted items, each represented by embedding vector $\mathbf{e}_i \in \mathbb{R}^D$. In classical recommendation systems, a user's entire history is compressed into a static vector via sum-pooling or average-pooling ($\mathbf{v}_U = \frac{1}{H}\sum \mathbf{e}_i$), which severely blurs diverse user interests.

A \textbf{Deep Interest Network (DIN)} introduces a candidate-dependent \textbf{Local Activation Unit} $a(\mathbf{e}_i, \mathbf{v}_A)$ that dynamically scales the contribution of each historical interaction relative to candidate $\mathbf{v}_A$:
\begin{enumerate}
    \item For each historical item $\mathbf{e}_i$, an explicit 4-way interaction vector is constructed:
    {\boldmath
    \begin{equation}
    \mathbf{I}(\mathbf{e}_i, \mathbf{v}_A) = \left[ \mathbf{e}_i; \; \mathbf{v}_A; \; \mathbf{e}_i \odot \mathbf{v}_A; \; \mathbf{e}_i - \mathbf{v}_A \right] \in \mathbb{R}^{4D}
    \end{equation}
    }
    \item The interaction is passed through an attention multilayer perceptron:
    {\boldmath
    \begin{equation}
    a(\mathbf{e}_i, \mathbf{v}_A) = \mathbf{w}_2^T \max\left( \mathbf{0}, \; W_1 \mathbf{I}(\mathbf{e}_i, \mathbf{v}_A) + \mathbf{b}_1 \right) + b_2
    \end{equation}
    }
    where $W_1 \in \mathbb{R}^{D_{\text{attn}} \times 4D}$, $\mathbf{b}_1 \in \mathbb{R}^{D_{\text{attn}}}$, $\mathbf{w}_2 \in \mathbb{R}^{D_{\text{attn}}}$, and $b_2 \in \mathbb{R}$.
    \item The activation coefficient $w_i \ge 0$ is obtained either through unnormalized positive activation:
    {\boldmath
    \begin{equation}
    w_i = \max\left( 0, \; a(\mathbf{e}_i, \mathbf{v}_A) \right)
    \end{equation}
    }
    or via normalized softmax attention.
    \item The dynamically conditioned \textbf{User Interest Representation} is:
    {\boldmath
    \begin{equation}
    \mathbf{v}_U(\mathbf{v}_A) = \sum_{i=1}^H w_i \mathbf{e}_i \in \mathbb{R}^D
    \end{equation}
    }
\end{enumerate}

\subsection{Expanded Non-Technical Definition (Intuition for Anyone)}
\textbf{Intuitive Conceptual Definition:}
\begin{enumerate}
    \item \textbf{The Fallacy of a Single User Profile}: A typical e-commerce shopper buys sneakers, cat food, computer keyboards, and gardening tools. If an AI averages all these together into a single summary vector, the vector points to a weird hybrid: a ``cat-food-gardening-shoe'' that matches nothing!
    \item \textbf{Candidate-Activated Spotlight}: When Amazon is deciding whether to show you a pair of running shoes, the cat food and gardening tools are completely irrelevant. DIN shines a spotlight that activates past shoe-related clicks while leaving everything else dim.
    \item \textbf{Unnormalized Attention Preserves Passion}: Traditional attention forces weights to sum to 1.0 (softmax). But if User Alice bought 50 pairs of shoes (an extreme enthusiast), and User Bob bought 1 pair (a casual buyer), their shoe-interest vectors should NOT have equal magnitude! DIN uses unnormalized attention so Alice's vector is $50\times$ stronger than Bob's.
    \item \textbf{Seamless Integration}: The customized interest vector $\mathbf{v}_U$ flows directly into standard click-through prediction heads.
\end{enumerate}

\subsection{Child-Friendly Mental Model (Physical Metaphor)}
Imagine having a magical trunk filled with toys:
\begin{itemize}
    \item \textbf{The Mixed Toybox}: Inside the trunk are toy cars, dinosaur figures, watercolor paints, and building blocks.
    \item \textbf{The Candidate Challenge}: A friend brings over a brand new remote-control monster truck ($\mathbf{v}_A$) and asks: ``Do you want to play?''
    \item \textbf{The Magnetic Magnet}: Instead of pulling out the paints or dinosaurs, your mind instantly connects with the toy cars and race tracks in your trunk ($\mathbf{v}_U$).
    \item \textbf{Immediate Enthusiasm}: Because you own 20 toy cars, your excitement level shines bright and you shout: ``YES!''
\end{itemize}

\section{Core Mathematical Equations and Definitions}

\subsection{Four-Way Interaction Tensor Construction}
{\boldmath
\begin{equation}
\mathbf{I}_i = \left[ \mathbf{e}_i; \; \mathbf{v}_A; \; \mathbf{e}_i \odot \mathbf{v}_A; \; \mathbf{e}_i - \mathbf{v}_A \right] \in \mathbb{R}^{4D}
\end{equation}
}
\begin{itemize}
    \item \textbf{Formal Definition}: The concatenated interaction feature combining individual representations, element-wise correlation, and vector difference.
    \item \textbf{What It Means}: Supplies the attention network with absolute positions, multiplicative alignment, and geometric distance simultaneously.
    \item \textbf{Why It Is Important}: Outperforms simple dot products by allowing the MLP to learn asymmetric and non-linear relevance boundaries.
\end{itemize}

\subsection{Local Activation Score Function}
{\boldmath
\begin{equation}
a_i = \mathbf{w}_2^T \operatorname{ReLU}\left( W_1 \mathbf{I}_i + \mathbf{b}_1 \right) + b_2
\end{equation}
}
\begin{itemize}
    \item \textbf{Formal Definition}: The feed-forward scalar regression evaluating candidate affinity for history item $i$.
    \item \textbf{What It Means}: Estimates the relevance scalar between the user's past interaction and the current candidate.
    \item \textbf{Why It Is Important}: Projects high-dimensional interaction manifolds into a scalar attention weight.
\end{itemize}

\subsection{Candidate-Conditioned User Interest Pooling}
{\boldmath
\begin{equation}
\mathbf{v}_U(\mathbf{v}_A) = \sum_{i=1}^H w_i \mathbf{e}_i
\end{equation}
}
\begin{itemize}
    \item \textbf{Formal Definition}: The weighted sum of historical item embeddings parameterized by candidate query $\mathbf{v}_A$.
    \item \textbf{What It Means}: Synthesizes an ad-hoc user representation custom-tailored to the specific item being evaluated.
    \item \textbf{Why It Is Important}: Resolves the fixed-capacity representation bottleneck of static user profile models.
\end{itemize}

\subsection{Total Attention Mass Metric}
{\boldmath
\begin{equation}
M_{\text{attn}}(\mathbf{v}_A) = \sum_{i=1}^H w_i
\end{equation}
}
\begin{itemize}
    \item \textbf{Formal Definition}: The unnormalized scalar sum of activation weights across the entire behavior history.
    \item \textbf{What It Means}: Directly quantifies user historical engagement intensity and interest volume within the candidate's domain.
    \item \textbf{Why It Is Important}: Preserves engagement frequency signals that are erased when using standard softmax normalization.
\end{itemize}

\section{Rigorous Mathematical Analysis Checklist}

\begin{itemize}
    \item \textbf{Notation}: $\mathbf{v}_A \in \mathbb{R}^D$ (candidate embedding), $\mathbf{e}_i \in \mathbb{R}^D$ (history embedding), $H$ (history length), $w_i$ (attention weight), $\mathbf{v}_U$ (user interest vector).
    \item \textbf{Epistemic Status}: Proposed by Zhou et al. (Alibaba, KDD 2018), deployed across massive commercial display advertising platforms.
    \item \textbf{Dependency Graph}: $(\mathbf{e}_i, \mathbf{v}_A) \to \mathbf{I}_i \to a_i \to w_i \to \mathbf{v}_U = \sum w_i \mathbf{e}_i \to \text{Top MLP Scoring}$.
    \item \textbf{Dimensions}: Item dimension $D$, interaction dimension $4D$, attention hidden dimension $D_{\text{attn}}$, history count $H$.
    \item \textbf{Regimes}: E-commerce CTR prediction with rich user interaction histories (clicks, add-to-carts, views).
    \item \textbf{Invariants}: If $\mathbf{e}_i \perp \mathbf{v}_A$ for all $i$ and biases are zero, $w_i \approx 0$ and $\mathbf{v}_U \approx \mathbf{0}$.
    \item \textbf{Isomorphisms}: DIN is isomorphic to cross-attention between a single target query token ($\mathbf{v}_A$) and a sequence of key/value tokens ($\mathbf{e}_i$) without causal masking.
\end{itemize}

\section{Theoretical Theorems and Formal Proofs}

\begin{theorem}[Candidate-Conditioned Dynamic Manifold Diversity]
Let $f(\mathbf{v}_A) \equiv \mathbf{v}_U(\mathbf{v}_A) = \sum_{i=1}^H w_i(\mathbf{v}_A) \mathbf{e}_i$ define the user interest mapping. If the historical set $\mathcal{H} = \{\mathbf{e}_1, \dots, \mathbf{e}_H\}$ contains at least two linearly independent vectors $\mathbf{e}_1 \not\propto \mathbf{e}_2$, then the image of $f$ has dimension strictly greater than 1:
{\boldmath
\begin{equation}
\dim\left( \operatorname{span}\left( \{ f(\mathbf{v}_A) \mid \mathbf{v}_A \in \mathbb{R}^D \} \right) \right) \ge 2
\end{equation}
}
In contrast, static pooling models $\mathbf{v}_U^{\text{static}} = \frac{1}{H}\sum_{i=1}^H \mathbf{e}_i$ have zero degrees of freedom ($\dim = 0$, a constant single point).
\end{theorem}

\begin{proof}
Consider two candidates $\mathbf{v}_A^{(1)} = \mathbf{e}_1$ and $\mathbf{v}_A^{(2)} = \mathbf{e}_2$.
By construction of the interaction vector:
{\boldmath
\begin{equation}
\mathbf{I}(\mathbf{e}_1, \mathbf{v}_A^{(1)}) = [\mathbf{e}_1; \mathbf{e}_1; \mathbf{e}_1 \odot \mathbf{e}_1; \mathbf{0}]
\end{equation}
}
Because the difference term $\mathbf{e}_1 - \mathbf{v}_A^{(1)} = \mathbf{0}$ and element-wise product is maximal, the trained activation unit assigns maximal weight $w_1(\mathbf{v}_A^{(1)}) = \alpha > 0$, while for orthogonal item $\mathbf{e}_2$, the difference $\|\mathbf{e}_2 - \mathbf{e}_1\|_2 > 0$ leads to $w_2(\mathbf{v}_A^{(1)}) = \epsilon \approx 0$.
Thus $f(\mathbf{v}_A^{(1)}) \approx \alpha \mathbf{e}_1$.
By identical symmetry, evaluating with candidate $\mathbf{v}_A^{(2)} = \mathbf{e}_2$ yields $w_2(\mathbf{v}_A^{(2)}) = \alpha > 0$ and $w_1(\mathbf{v}_A^{(2)}) \approx 0$, so $f(\mathbf{v}_A^{(2)}) \approx \alpha \mathbf{e}_2$.
Since $\mathbf{e}_1$ and $\mathbf{e}_2$ are linearly independent, the span of $\{f(\mathbf{v}_A^{(1)}), f(\mathbf{v}_A^{(2)})\}$ has dimension 2:
{\boldmath
\begin{equation}
\dim\left( \operatorname{span}\left( \{ f(\mathbf{v}_A) \} \right) \right) \ge 2
\end{equation}
}
For static pooling, $\mathbf{v}_U^{\text{static}}$ is a constant vector independent of $\mathbf{v}_A$, with zero variance across candidates.
\end{proof}

\begin{theorem}[Preservation of Intensity Order under Unnormalized Summation]
Let User 1 have history $\mathcal{H}_1 = \{\mathbf{e}\}$ (1 interaction) and User 2 have history $\mathcal{H}_2 = \{\mathbf{e}, \mathbf{e}, \dots, \mathbf{e}\}$ ($N$ identical interactions). Under unnormalized DIN attention with positive activation $w(\mathbf{e}, \mathbf{v}_A) = c > 0$:
{\boldmath
\begin{equation}
\|\mathbf{v}_{U, 2}(\mathbf{v}_A)\|_2 = N \cdot \|\mathbf{v}_{U, 1}(\mathbf{v}_A)\|_2
\end{equation}
}
Under standard softmax normalized attention, $\|\mathbf{v}_{U, 2}(\mathbf{v}_A)\|_2 = \|\mathbf{v}_{U, 1}(\mathbf{v}_A)\|_2$, destroying the frequency signal.
\end{theorem}

\begin{proof}
For User 1, the unnormalized sum is:
{\boldmath
\begin{equation}
\mathbf{v}_{U, 1} = w(\mathbf{e}, \mathbf{v}_A) \mathbf{e} = c \mathbf{e} \implies \|\mathbf{v}_{U, 1}\|_2 = c \|\mathbf{e}\|_2
\end{equation}
}
For User 2 with $N$ identical items, the unnormalized sum is:
{\boldmath
\begin{equation}
\mathbf{v}_{U, 2} = \sum_{i=1}^N w(\mathbf{e}, \mathbf{v}_A) \mathbf{e} = \sum_{i=1}^N c \mathbf{e} = N c \mathbf{e} \implies \|\mathbf{v}_{U, 2}\|_2 = N c \|\mathbf{e}\|_2 = N \|\mathbf{v}_{U, 1}\|_2
\end{equation}
}
Now consider softmax normalization.
For User 2, the softmax weight for each item is:
{\boldmath
\begin{equation}
\widetilde{w}_i = \frac{\exp(s)}{\sum_{j=1}^N \exp(s)} = \frac{\exp(s)}{N \exp(s)} = \frac{1}{N}
\end{equation}
}
The resulting normalized interest vector is:
{\boldmath
\begin{equation}
\widetilde{\mathbf{v}}_{U, 2} = \sum_{i=1}^N \frac{1}{N} \mathbf{e} = \mathbf{e} = \widetilde{\mathbf{v}}_{U, 1}
\end{equation}
}
Thus softmax normalization yields identical representations for 1 interaction and $N$ interactions, completely erasing behavioral frequency information.
\end{proof}

\section{Foundational Architectural Questions and Epistemic Meta-Questions}

\subsection{Architectural Question 1: Unnormalized vs Softmax Normalized Attention}
\textbf{Question:} Why did Alibaba's DIN explicitly abandon standard softmax attention in favor of unnormalized attention?\\
\textbf{Answer:} Softmax enforces $\sum w_i = 1$, which acts as a mathematical average. If a user is generally interested in sports (100 clicks) and casually clicked a shoe once (1 click), when evaluating sports ads, softmax assigns similar normalized weights regardless of whether the user clicked 5 times or 500 times. Unnormalized attention preserves total attention mass $\sum w_i$, providing a direct proxy for user interest volume and engagement strength.

\textbf{Meta-Question:} Doesn't unnormalized attention cause activation explosion on very long histories?\\
\textbf{Answer:} Yes. On histories with $H > 1000$, unnormalized attention can produce large output magnitudes that saturate downstream sigmoid layers. This is controlled via layer normalization, $L_2$ regularization, or logarithmic scaling of the aggregated vector.

\subsection{Architectural Question 2: Lifelong Behavior Modeling (SIM and ETA)}
\textbf{Question:} How is target attention scaled when users have tens of thousands of historical clicks ($H > 50,000$)?\\
\textbf{Answer:} Evaluating the 4-way interaction MLP over 50,000 items per candidate during live serving violates strict latency budgets ($<20$ ms). Modern architectures (Search-based Interest Model, SIM, and ETA) use a two-stage approach: (1) General Search Unit (GSU) retrieves the top 100 most relevant items using category filters or vector indexes, and (2) Exact DIN attention is applied only over those top 100 items.

\textbf{Meta-Question:} How does DIN compare to Deep Interest Evolution Network (DIEN)?\\
\textbf{Answer:} DIEN extends DIN by adding an auxiliary loss and an Attention-over-GRU (AUGRU) to model sequential interest transitions over time, capturing how a user's interest evolved from casual browsing to committed purchase.

\section{Algorithmic Implementation Specification}

\begin{algorithm}[H]
\caption{Deep Interest Network (DIN) Target-Dependent Attention Forward Pass}
\begin{algorithmic}[1]
\Require Candidate embedding $\mathbf{v}_A \in \mathbb{R}^D$, history embeddings $E \in \mathbb{R}^{H \times D}$, attention weights $W_1 \in \mathbb{R}^{D_{\text{attn}} \times 4D}$, $\mathbf{b}_1 \in \mathbb{R}^{D_{\text{attn}}}$, $\mathbf{w}_2 \in \mathbb{R}^{D_{\text{attn}}}$, bias $b_2 \in \mathbb{R}$, softmax flag $flag_{\text{norm}}$
\Ensure User interest vector $\mathbf{v}_U \in \mathbb{R}^D$, attention weights $\mathbf{w} \in \mathbb{R}^H$, total mass $M$, dominant item index $i^*$
\State Initialize raw scores $\mathbf{s} \in \mathbb{R}^H$
\For{$i = 0$ \textbf{to} $H-1$} \Comment{Evaluate Local Activation Units}
    \State $\mathbf{e}_i \gets E[i]$
    \State $\mathbf{I} \gets [\mathbf{e}_i; \; \mathbf{v}_A; \; \mathbf{e}_i \odot \mathbf{v}_A; \; \mathbf{e}_i - \mathbf{v}_A] \in \mathbb{R}^{4D}$
    \State Initialize hidden vector $\mathbf{h} \in \mathbb{R}^{D_{\text{attn}}}$
    \For{$j = 0$ \textbf{to} $D_{\text{attn}}-1$}
        \State $val \gets \sum_{k=0}^{4D-1} W_1[j][k] \cdot \mathbf{I}[k] + \mathbf{b}_1[j]$
        \State $\mathbf{h}[j] \gets \max(0.0, val)$
    \EndFor
    \State $\mathbf{s}[i] \gets \sum_{j=0}^{D_{\text{attn}}-1} \mathbf{w}_2[j] \cdot \mathbf{h}[j] + b_2$
\EndFor
\State Initialize $\mathbf{w} \in \mathbb{R}^H$
\If{$flag_{\text{norm}}$ is true} \Comment{Softmax Normalization}
    \State $s_{\max} \gets \max_i \mathbf{s}[i]$
    \State $sum\_exp \gets \sum_{i=0}^{H-1} \exp(\mathbf{s}[i] - s_{\max})$
    \For{$i = 0$ \textbf{to} $H-1$}
        \State $\mathbf{w}[i] \gets \exp(\mathbf{s}[i] - s_{\max}) / sum\_exp$
    \EndFor
\Else \Comment{Unnormalized Positive Activation}
    \For{$i = 0$ \textbf{to} $H-1$}
        \State $\mathbf{w}[i] \gets \max(0.0, \mathbf{s}[i])$
    \EndFor
\EndIf
\State Initialize $\mathbf{v}_U \in \mathbb{R}^D$ to zeros
\For{$i = 0$ \textbf{to} $H-1$} \Comment{Target-Dependent Weighted Aggregation}
    \State $weight \gets \mathbf{w}[i]$
    \For{$d = 0$ \textbf{to} $D-1$}
        \State $\mathbf{v}_U[d] \gets \mathbf{v}_U[d] + weight \cdot E[i][d]$
    \EndFor
\EndFor
\State $M \gets \sum_{i=0}^{H-1} \mathbf{w}[i], \quad i^* \gets \arg\max_{0 \le i < H} \mathbf{w}[i]$
\State \Return $(\mathbf{v}_U, \mathbf{w}, M, i^*)$
\end{algorithmic}
\end{algorithm}

\section{Big-$\Theta$ Complexity Analysis}

\begin{itemize}
    \item \textbf{Activation Unit Time Complexity}: $\Theta(H \cdot (D_{\text{attn}} \cdot D))$ operations across $H$ historical interactions.
    \item \textbf{History Aggregation Complexity}: $\Theta(H \cdot D)$ operations.
    \item \textbf{Space Complexity}: $\Theta(H + D + D_{\text{attn}})$ for activation scores and interaction tensors.
    \item \textbf{Parallel Depth}: $\mathcal{D} = \Theta(\log D + \log D_{\text{attn}} + \log H)$.
\end{itemize}

\section{Single Canonical Repository Reference}

The complete deterministic scalar implementation, verification suite, and unit test benchmarks are published at:
\begin{center}
\url{https://github.com/Chief-Strategist-J/llm-obs-infra}
\end{center}

\end{document}
